% paper/sections/12f_interface_crossing.tex
%
% §12.5  高密度比対応：一括解法の限界
% ═══════════════════════════════════════════════════════════════════════════════
\subsection{高密度比対応：変密度一括解法の適用限界}
\label{sec:interface_crossing}
% ═══════════════════════════════════════════════════════════════════════════════

§~\ref{sec:split_ppe_framework} で分相 PPE 解法を主要 PPE 戦略として定式化し，
§~\ref{sec:verify_split_ppe_dc} でその密度比非依存な精度（$\Ord{h^{7.0}}$）を実証した．
本節では，変密度 smoothed Heaviside 一括解法を再試験し，
旧版で報告した「$\rho_l/\rho_g \geq 10$ で DC が発散する」という結果が
再現するかを検証する．

% ─────────────────────────────────────────────────────────────────────────────
\subsubsection{界面型密度場での DC 収束特性}
\label{sec:varrho_dc_convergence}
% ─────────────────────────────────────────────────────────────────────────────

\noindent\textbf{目的：}
変密度 smoothed Heaviside 一括解法の欠陥補正（DC）が
界面型密度場でどの程度まで残差を下げられるかを定量化する．

\noindent\textbf{評価対象：}
変密度 PPE 欠陥補正（DC $k = 3$）の密度比依存性．

\noindent\textbf{手順とパラメタ：}
\begin{itemize}[noitemsep]
  \item 製造解 $p^* = \sin\pi x\sin\pi y$，ディリクレ BC
  \item 密度場 $\rho = \rho_g + (\rho_l - \rho_g)\,H_\varepsilon(\phi)$（$\varepsilon = 1.5h$，円形界面 $R = 0.25$）
  \item 高次演算子
        $L_H(p) = \rho^{-1}\nabla^2 p - (\bnabla\rho\cdot\bnabla p)/\rho^2$
        と FD 変係数前処理を用いた DC $k = 3$
  \item $N = 64$，密度比 $\rho_l/\rho_g = 2, 5, 10, 20, 50, 100$
\end{itemize}

\noindent\textbf{結果：}

\begin{table}[htbp]
\centering
\caption{変密度 PPE 欠陥補正の密度比依存性（$N = 64$，$k = 3$，界面型密度場）}
\label{tab:varrho_breakdown}
\renewcommand{\arraystretch}{1.3}
\begin{tabular}{rcccc}
\toprule
$\rho_l/\rho_g$ & $L^\infty$ 誤差（液相内部） & $\|r_0\|_\infty$ & $\|r_3\|_\infty$ & 判定 \\
\midrule
   2 & $8.4 \times 10^{-8}$ & $4.5 \times 10^{1}$ & $1.4 \times 10^{-2}$ & 収束 \\
   5 & $5.8 \times 10^{-7}$ & $1.4 \times 10^{2}$ & $5.4 \times 10^{-2}$ & 収束 \\
  10 & $2.9 \times 10^{-7}$ & $3.1 \times 10^{2}$ & $8.8 \times 10^{-2}$ & 収束 \\
  20 & $1.4 \times 10^{-6}$ & $6.5 \times 10^{2}$ & $1.1 \times 10^{-1}$ & 収束 \\
  50 & $5.0 \times 10^{-6}$ & $1.7 \times 10^{3}$ & $1.5 \times 10^{-1}$ & 収束 \\
 100 & $9.1 \times 10^{-6}$ & $3.4 \times 10^{3}$ & $2.8 \times 10^{-1}$ & 収束 \\
\bottomrule
\end{tabular}
\end{table}

\noindent\textbf{結論：}
旧版の「$\rho_l/\rho_g \geq 10$ で DC が発散する」という結果は，
ゼロベースで作り直した再試験では再現しなかった．
DC $k=3$ は $\rho_l/\rho_g = 100$ まで残差を $10^{-1}$ 程度まで低下させ，
液相内部誤差も $10^{-5}$ 以下に留まった．
したがって，本テストから主張できる限界は
\textbf{密度比 10 以上での発散}ではなく，
\textbf{界面型密度場における漸近格子収束の喪失}である．

% ─────────────────────────────────────────────────────────────────────────────
\subsubsection{格子収束性の密度比依存}
\label{sec:varrho_grid_convergence}
% ─────────────────────────────────────────────────────────────────────────────

\noindent\textbf{目的：}
DC が残差を下げる低--中密度比（$\rho_l/\rho_g = 2, 5$）について格子収束次数を計測し，
CCD の設計精度 $\Ord{h^6}$ が界面型密度場で回復可能かを検証する．

\noindent\textbf{評価対象：}
格子収束次数の密度比依存性（変密度一括 PPE + CCD 離散化）．

\noindent\textbf{手順とパラメタ：}
\begin{itemize}[noitemsep]
  \item 前節と同一の製造解（$p^* = \sin\pi x\sin\pi y$，界面型密度場）
  \item $\rho_l/\rho_g = 2, 5$，$N = 16, 32, 64, 128$
  \item 静止液滴密度比スイープ：$\rho_l/\rho_g = 2, 3, 5, 10$，$N = 64$，$\mathit{We} = 10$，50 ステップ
\end{itemize}

\noindent\textbf{結果：}

$\rho_l/\rho_g = 2$ では
$N=16 \to 32$ で $5.53$ 次，$32 \to 64$ で $4.38$ 次を示したが，
$64 \to 128$ では $-0.64$ 次となり誤差がプラトー化した．
$\rho_l/\rho_g = 5$ でも同様に，
$5.62$ 次，$2.26$ 次，$-0.17$ 次と高解像度側で漸近収束が失われた．
界面幅 $\varepsilon = 1.5h$ が格子細分化とともに縮小するため，
$\bnabla\rho$ の急峻化と界面横断ステンシルの汚染が高次収束を阻害する．

図~\ref{fig:ch12_density_fields} に静止液滴の圧力場と寄生電流を
$\rho_l/\rho_g = 2, 3, 5, 10$ で比較する．
$\rho_l/\rho_g \leq 10$ では Laplace 圧が維持され，
寄生流れは有界に保たれる．
一方で，この安定性は上記 MMS で確認した格子収束の喪失を否定しない．

\begin{figure}[htbp]
\centering
\includegraphics[width=\linewidth]{figures/ch12_density_fields}
\caption{密度比スイープ：静止液滴の圧力場（上段）と寄生電流（下段）．
$N = 64$，$\mathit{We} = 10$，50 ステップ．
$\rho_l/\rho_g = 2, 3, 5, 10$ の全ケースで安定．}
\label{fig:ch12_density_fields}
\end{figure}

\begin{tcolorbox}[resultbox, title={Smoothed Heaviside 一括 DC の再試験結果}]
\begin{enumerate}[noitemsep]
  \item DC $k{=}3$ の発散閾値は本再試験では再現されず，
        $\rho_l/\rho_g = 100$ まで残差は低下した
  \item ただし界面型密度場では高解像度側で格子収束がプラトー化し，
        CCD の設計精度 $\Ord{h^6}$ は回復しない
  \item 分相 PPE + DC $k{=}3$（§~\ref{sec:split_ppe_framework}）は
        各相内部で $\Ord{h^{7.0}}$ を達成するため，
        高密度比・高精度計算では引き続き主要戦略となる
\end{enumerate}
ただし分相 PPE は\textbf{相ごとに定密度 Poisson 方程式を個別に解く}手法であり，
smoothed Heaviside の変密度演算子を定密度に置換するだけでは
界面での $(\bnabla\rho \cdot \bnabla p)/\rho^2$ 項の欠落により
速度発散誤差が蓄積する．
実装には GFM 等の界面ジャンプ処理が不可欠である．
\end{tcolorbox}

\noindent\textbf{結論：}
変密度一括 PPE の主要な制約は，今回の再試験では
「密度比 10 以上での DC 発散」ではなく，
界面型密度場における漸近格子収束の喪失として確認された．
この制約は界面を横切る CCD/FD ステンシルに由来し，
高密度比の実用計算では分相 PPE と界面ジャンプ処理を用いるべきである．
